\begin{register}{H}{dcsr}{0x7B0}\label{regdesc:LPDCSR}
\regfield{xdebugver}{4}{28}{0100}%
\regfield{(reserved)}{12}{16}{000000000000}%
\regfield{ebreakm}{1}{15}{0}%
\regfield{(reserved)}{2}{13}{00}%
\regfield{ebreaku}{1}{12}{0}%
\regfield{(reserved)}{3}{9}{000}%
\regfield{cause}{3}{6}{000}%
\regfield{(reserved)}{3}{3}{000}%
\regfield{step}{1}{2}{0}%
\regfield{prv}{2}{0}{11}%
\reglabel{Reset}\regnewline%
\begin{regdesc}\begin{reglist}
\label{xdebugver}\item [xdebugver] Represents the debug version.\\
4: External debug support exists\\
(RO)
\label{ebreakm}\item [ebreakm] Configures execution of the EBREAK instruction in machine mode.\\
0: Trigger an exception with mcause = 3\\
1: Enter debug mode\\
(R/W)
\label{ebreaku}\item [ebreaku] Configures execution of the EBREAK instruction in user mode.\\
0: Trigger an exception with mcause = 3 as described in privileged mode\\
1: Enter debug mode\\
(R/W)
\label{cause}\item [cause] Represents the reason why debug mode was entered. When there are multiple reasons to enter debug mode in a single cycle, the cause with the highest priority number is the one written.\\
1: An EBREAK instruction was executed. (priority 3)\\
2: The Trigger Module caused a halt. (priority 4)\\
3: haltreq was set. (priority 2)\\
4: The CPU core single stepped because step was set. (priority 1)\\
Other values: reserved for future use\\
(RO)
\label{step}\item [step] When set and not in Debug Mode, the core will only execute a single instruction and then enter Debug Mode. \\
If the instruction does not complete due to an exception, the core will immediately enter Debug Mode before executing the trap handler, with appropriate exception registers set.\\
Setting this bit does not mask interrupts. This is a deviation from the RISC-V External Debug Support Specification Version 0.13.\\
(R/W)
\label{prv}\item [prv] Contains the privilege level the core is operating in when debug mode is entered. A debugger can change this value to change the core's privilege level when exiting debug mode. Only \textbf{0x3} (machine mode) and \textbf{0x0} (user mode) are supported. (RO)
\end{reglist}\end{regdesc}
\end{register}

%%%%%%%%%
\begin{register}{H}{dpc}{0x7B1}\label{regdesc:LPDPC}
\regfield{dpc}{32}{0}{00000000000000000000000000000000}%
\reglabel{Reset}\regnewline%
\begin{regdesc}\begin{reglist}
\label{dpcr}\item [dpc] Upon entry to debug mode, dpc is written with the address of the next instruction that will be executed. When resuming, the CPU core's PC is updated to the address stored in dpc. In debug mode, dpc can be modified. This field can be accessed in debug mode. (R/W)
\end{reglist}\end{regdesc}
\end{register}


%%%%%%%%%
\begin{register}{H}{dscratch0}{0x7B2}\label{regdesc:LPDSCRATCH0}
\regfield{dscratch0}{32}{0}{00000000000000000000000000000000}%
\reglabel{Reset}\regnewline%
\begin{regdesc}\begin{reglist}
\label{dscratch0}\item [dscratch0] Used by the debug module internally. (R/W)
\end{reglist}\end{regdesc}
\end{register}

%%%%%%%%%
\begin{register}{H}{dscratch1}{0x7B3}\label{regdesc:LPDSCRATCH1}
\regfield{dscratch1}{32}{0}{00000000000000000000000000000000}%
\reglabel{Reset}\regnewline%
\begin{regdesc}\begin{reglist}
\label{dscratch1}\item [dscratch1] Used by the debug module internally.(R/W)
\end{reglist}\end{regdesc}
\end{register}
